% =====================================================================
% Section label naming convention (kebab-case, parent prefix per level).
% Cross-references use section labels; update this map when sections move.
%
%   Section 1  Opening Movement                          sec:opening
%   Section 2  The Frame as Method                       sec:frame-as-method
%     2.1  Two kinds of absence                          sec:frame-as-method:two-kinds-of-absence
%     2.2  The empty chairs in BSA/AML monitoring        sec:frame-as-method:bsa-aml-empty-chairs
%     2.3  Walking three controls through the frame     sec:frame-as-method:walking-three-controls
%     2.4  What the frame produced                       sec:frame-as-method:what-the-frame-produced
%   Section 3  Verification Regimes as Taxonomic Move    sec:verification-regimes
%     3.1  The structural line                           sec:verification-regimes:structural-line
%     3.2  Adverse action notices as the canonical case  sec:verification-regimes:adverse-action-notices
%     3.3  The taxonomic move                            sec:verification-regimes:taxonomic-move
%     3.4  What the categoricals foreclose, and not      sec:verification-regimes:foreclosures
%     3.5  A turn: pricing as produced-absence case      sec:verification-regimes:pricing-turn
%     3.6  The taxonomic move, summarized                sec:verification-regimes:summary
%   Section 4  Diagnostic Techniques                     sec:diagnostics
%     4.1  Silence-manufacture as structural pattern     sec:diagnostics:silence-manufacture
%       4.1.1  Definition                                sec:diagnostics:silence-manufacture:definition
%       4.1.2  Relationship to Goodhart's law            sec:diagnostics:silence-manufacture:goodhart
%       4.1.3  Mapping the three earned positions        sec:diagnostics:silence-manufacture:mapping-positions
%       4.1.4  What silence-manufacture predicts         sec:diagnostics:silence-manufacture:predictions
%     4.2  Three diagnostic techniques                   sec:diagnostics:three-techniques
%       4.2.1  Technique 1: Empty-chair enumeration      sec:diagnostics:enumeration
%       4.2.2  Technique 2: Pages missing from ledger    sec:diagnostics:pages-missing
%       4.2.3  Technique 3: Produced-silence pair        sec:diagnostics:question-pair
%     4.3  From orientation to evidence request          sec:diagnostics:evidence-request
%     4.4  Coordination consequence                      sec:diagnostics:compounding
%   Section 5  Provisional Evidentiary Capabilities      sec:primitives
%     5.1  Revision-aware temporal capture               sec:primitives:temporal-capture
%     5.2  Evidence binding and exportability            sec:primitives:evidence-binding
%     5.3  Non-collapsing uncertainty representation     sec:primitives:confidence-characterization
%     5.4  Technical-change monitoring                   sec:primitives:technical-change
%     5.5  Institutional-conformance monitoring          sec:primitives:institutional-conformance
%     5.6  What the capabilities contribute              sec:primitives:compound
%   Section 6  Honest Uncertainties                      sec:uncertainties
%     6.1  Uncertainties about the frame                 sec:uncertainties:frame
%     6.2  Uncertainties about the categoricals          sec:uncertainties:categoricals
%     6.3  Uncertainties about the diagnostic techniques sec:uncertainties:diagnostics
%     6.4  Architectural boundaries
%     6.5  The meta-uncertainty: theory of effect        sec:uncertainties:meta
%     6.6  What this section earns                       sec:uncertainties:summary
%   Section 7  Conclusion                                sec:implications
% =====================================================================

\documentclass[11pt]{article}

\usepackage[T1]{fontenc}
\usepackage[english]{babel}
\usepackage[margin=1in]{geometry}
\usepackage{microtype}
\usepackage{csquotes}
\usepackage[backend=biber,style=authoryear,maxcitenames=2,giveninits=true,uniquename=false]{biblatex}
\addbibresource{references.bib}
\usepackage{hyperref}
\hypersetup{
  colorlinks=true,
  linkcolor=black,
  citecolor=black,
  urlcolor=blue,
  pdftitle={Architectures of Absence: AI Governance under the FS AI RMF},
  pdfauthor={Tony Mason}
}

\title{Architectures of Absence:\\AI Governance under the FS~AI~RMF}
\author{Tony Mason\\
  \texttt{fsgeek@gatech.edu}}
\date{July 23, 2026}

\begin{document}

\maketitle

\begin{abstract}
The Financial Services AI Risk Management Framework (FS~AI~RMF) gives community
banks 230 control objectives but not a common language for deciding what their
artifacts establish. Banks may also depend on vendor systems they are responsible
for governing but cannot unilaterally change. This position paper offers a
translation ontology for coordination among banks, examiners and regulators,
and vendors and architects, with affected parties as its normative center.
\emph{Empty-chair representation} asks whose interest a control protects when
that party is absent from an AI-supported decision.

The ontology follows an affected interest through control objective,
evidentiary object, available artifact, claimed inference and binding, and the
actor able to close, disclose, or contest a gap. Decision records, feature
attributions, local approximations, generated rationales, human review records,
and approvals expose different objects; none inherits another's warrant merely
by being called an explanation. \emph{Silence-manufacture} names the diagnostic
sequence in which unavailable contradictory evidence helps an unsupported
substitution persist.

Three two-layer diagnostics first orient inquiry and then structure any
authorized evidence request: empty-chair enumeration, pages missing from the
ledger relative to a declared attestation scope, and the produced-silence
question pair. They motivate five provisional evidentiary capabilities:
revision-aware temporal capture; evidence binding and exportability;
non-collapsing uncertainty representation; technical-change monitoring; and
institutional-conformance monitoring. Worked applications to selected
FS~AI~RMF controls illustrate the vocabulary rather than validate it across the
framework. The paper supplies neither an implementation validation nor new
examination authority; it identifies where current claims, technical
capabilities, institutional obligations, and policy choices should meet.
\end{abstract}

% Section 1: Opening Movement
%
% Regulator-facing statement of the problem, lens, and evidentiary taxonomy.

\section{Opening Movement}
\label{sec:opening}

The Financial Services AI Risk Management Framework (FS~AI~RMF), developed by the Cyber Risk Institute with the Financial Services Sector Coordinating Council and the US Treasury and released in February 2026, specifies 230 control objectives spanning governance, data, model development, validation, monitoring, third-party risk, and consumer protection \parencite{treasuryFSAIRMFRelease2026}. Read as a checklist, the framework is operationally daunting and conceptually opaque: community banks still face the question of what \emph{commensurate} compliance means, and the document itself does not answer it. The FS~AI~RMF, unlike the supervisory model-risk guidance issued the same year (SR Letter~26-2), has no asset-size threshold; it is designed to apply to all financial institutions regardless of size, type, complexity, or criticality. Read as a structure, however, the framework admits a generative organizing lens that makes interpretation tractable and gives operationalization a structure to build on: a wide range of control objectives can be understood as protecting an interest of a party who is not in the room when an AI system makes or supports a financial decision. We develop empty-chair representation as an interpretive lens, demonstrated through worked examples across control families rather than asserted as a coding of all 230 objectives; whether the reading holds framework-wide is an empirical question we frame, but do not settle, here (a systematic content-analytic coding with trained raters is future work).

Recent commentary on the FS~AI~RMF describes it as a \emph{systems blueprint} rather than a checklist, and urges institutions to map controls to architecture, ownership, and evidence rather than to documentation alone~\parencite{mushahwar2026}. We extend that move one step further: the blueprint has a normative structure, and that structure is \emph{empty-chair representation}: the design of AI governance such that specific absent parties are structurally represented in the architecture of decision-making, not merely in its documentation.

The party in the empty chair varies by control. In some, it is the loan applicant who will be denied without recourse. In others, the depositor whose funds are deployed in ways they did not consent to. In others, the future compliance officer who will inherit responsibility for decisions made under regimes they did not construct. In others, the examiner applying updated model-risk guidance that had not yet been issued when the system was deployed. In others, the institution itself in its future state, after personnel turnover and ownership changes. The empty chair is plural, and the chairs do not always agree. The proposed architectural function of regulatory governance is to make these absent interests inspectable by the parties present at the moment of decision and by those who later review it.

This reframe matters operationally. By \emph{verification regime} we mean the set of technical, procedural, and evidentiary mechanisms by which a party not present at the moment of decision can later test a material claim about that decision, rather than accept the claim from the outcome or from an institution's assertion. A control \emph{represents} an absent party when it induces observable, auditable constraints on system behavior or evidence production that a later examiner can use to assess whether that party's interests were respected. Representation in this sense is architectural; documentation can participate in it, but an assertion of representation is not evidence that representation occurred.

The object being verified must be named. A decision record may preserve the \emph{inputs and context} available at decision time. SHAP can characterize \emph{functional feature dependence} under a specified value function and background construction; LIME can provide a \emph{local behavioral approximation} under a specified neighborhood and surrogate \parencite{lundberg2017,ribeiro2016}. Attention and activation analyses expose selected \emph{internal computational states}, although attention weight alone does not establish feature importance \parencite{jain2019}. Chain-of-thought and other natural-language accounts are \emph{generated rationales}. A contemporaneous reviewer record is evidence about a \emph{human deliberative record}: what was presented, consulted, recorded, and decided, not transparent access to every cognitive process. A signature or approval log records \emph{institutional authorization and review}. These artifacts can each be useful. They are not interchangeable, and an artifact may faithfully establish one object without establishing another.

This distinction changes the examination question. The question is not simply whether the institution has produced an explanation. It is \emph{what evidentiary object does this artifact expose, what conclusion is the recipient being asked to draw from it, and what binds the artifact to that conclusion?} A local approximation may illuminate behavior without recording deliberation; an understandable rationale may communicate a decision without faithfully characterizing functional dependence; a signature may establish authorization without establishing the truth of the material reviewed. Verification fails when the claimed conclusion requires a different object from the one the architecture exposes.

Examiners reviewing AI deployments at community banks face a problem the framework does not solve for them: 230 control objectives do not themselves specify how documentary evidence warrants a conclusion about deployed behavior. With an empty-chair frame, examination becomes a structured inquiry: \emph{which absent parties does this AI system affect, what artifact represents their interests, and what conclusion can that artifact support?} The frame supplements the checklist with an evidentiary question.

Two positions follow. First, \emph{an explanation artifact is not, by itself, a verification regime for a different evidentiary object}. Post-hoc methods have legitimate uses, including model diagnosis, communication, and exploration; their limits depend on their target and construction \parencite{oh2024}. But in adversarial review the recipient cannot assume that an understandable artifact establishes a different proposition without an inspectable binding between them \parencite{bordt2022}. SHAP and LIME can be manipulated in ways that conceal objectionable model behavior \parencite{slack2020}; generated chain-of-thought can be unfaithful to the computation that produced an answer \parencite{lanham2023,turpin2023}. These findings concern different artifacts and mechanisms. Their common governance lesson is narrower: the institution must establish the warrant for the inference it asks the examiner or affected party to draw.

Second, \emph{review of an explanation verifies that the explanation was reviewed; it does not by itself verify a decision process that the explanation does not faithfully represent}. A signature can establish authorization, responsibility, and completion of a specified procedure. Those are real governance functions. If the institution also claims that the signature establishes the substantive basis of an AI recommendation, it must show what the reviewer could inspect, what independent judgment the procedure required, and how the reviewed artifact bears on that substantive claim. Human presence in a workflow and verification of a model-supported decision are not equivalent propositions.

These positions do not counsel against AI in banking. They distinguish uses by the evidentiary claim the architecture must support. Pattern recognition, anomaly detection, and consistency monitoring can surface material for investigation while leaving the disposition and its documented basis with an authorized human. Decision-recommendation systems may demand evidence about model behavior, data, policy fit, human review, and the relation among them. The difference is not captured by a generic ``human in the loop'' label or a generic explanation requirement; it depends on which party decides, which artifact is exposed, and which conclusion that artifact can support.

A note on what the empty-chair frame is doing structurally, before the body sections develop it. Absence in regulated decision contexts is not always a static gap. An interface may make challenge costly; a documentation regime may preserve only later justification; a metric may omit the evidence needed to test the conclusion drawn from it. The frame therefore asks not only \emph{who is absent} but \emph{what access does the architecture remove, what artifact replaces it, and what inference is drawn from the replacement?} Section~\ref{sec:diagnostics} develops this recurring inference gap as silence-manufacture. The claim is conditional: silence is not evidence of absence when the system being evaluated materially contributes to that silence.

The remainder of this paper demonstrates the lens against specific FS~AI~RMF control objectives, develops three diagnostic questions, and derives four provisional architectural requirements at the property level. Its test is practical: whether the lens helps a community bank or examiner distinguish the evidentiary object an artifact exposes from the broader conclusion the institution asks that artifact to support.

% Section 2: The Frame as Method
%
% Status: first-pass draft integrating produced-silence distinction. Will be
% reworked.
%
% Citations needed: SAR data flagging biases, reviewer self-report literature
% (Nisbett \& Wilson 1977 plus operational decision-making), disparate-impact in
% BSA/AML monitoring, cross-institution displacement effects in money laundering.

\section{The Frame as Method}
\label{sec:frame-as-method}

The empty-chair frame's value depends on its generativity. If naming absent-party representation as the organizing principle behind the FS~AI~RMF produces only a more sympathetic vocabulary for compliance theater, the frame has earned nothing. If it produces interpretive leverage on specific controls --- vocabulary for talking about what a control is for, criteria for assessing whether an implementation represents the parties the control was written to protect, diagnostic moves available to examiners reviewing the implementation --- the frame has earned its position as an organizing principle. This section demonstrates the latter through a worked example in BSA/AML transaction monitoring, and in the course of doing so refines the frame itself by distinguishing two kinds of absence.

We choose BSA/AML rather than lending or underwriting deliberately. Lending is where the empty-chair frame is most obvious: the loan applicant is concrete, the harm is concrete, fair-lending controls have decades of regulatory history to draw on, and the absent party's interest is well-rehearsed. The frame applied to lending tends to confirm what readers already believe, which is a poor test of whether the frame is doing work or merely renaming work that was already being done. BSA/AML is harder. The empty chairs are plural and conflicting in ways lending is not. The institution's interest is more aligned with the regulator's than in lending. The customer's interest is ambiguous: protection from fraud and protection from false flagging point in opposite operational directions. If the frame produces clarity here, where the chairs disagree and the controls do not align cleanly with any single absent party, the frame is doing more than renaming.

\subsection{Two kinds of absence}
\label{sec:frame-as-method:two-kinds-of-absence}

Before enumerating the empty chairs in BSA/AML, we refine the frame. The opening movement treated absence as a single category: parties not in the room when a decision is made. In practice, absence has two structurally distinct forms, and the diagnostic moves available to address them differ.

\emph{Structural absence} describes parties whose absence is given by the situation: parties who do not yet exist, who have no embodied agent, or whose absence is a designed feature of the regulatory regime. The future compliance officer, the not-yet-victimized customer, the financial system as such, the criminal we are designing to exclude: these are absent because the architecture of the situation places them outside the decision context. The frame's diagnostic move on structural absence is \emph{how does the architecture make this party audible to the parties who are present?}

\emph{Produced absence} describes parties whose absence is created by the architecture itself, then read by the architecture as evidence of correctness. The customer who would have appealed if appeal were not prohibitively expensive. The borrower who stopped applying after repeated denial and whose absence from the application data is read as evidence of self-selection out of credit demand. The reviewer whose actual reasoning is structurally erased by a documentation regime that captures only post-hoc justification, where the absence of contradicting evidence in the record is read as evidence the documentation is faithful. In each case, the architecture has costs or design choices that suppress a voice, and the architecture then interprets the suppressed voice's silence as endorsement of the decision that suppressed it.

The distinction matters because produced absence is invisible to the empty-chair frame as initially stated. \emph{Make the absent party audible} is the wrong operation when the architecture is producing the silence: the chair appears empty because the architecture removed it, and asking how to give voice to the chair while leaving the removal-mechanism intact will fail. The diagnostic move on produced absence is \emph{what is the architecture doing that creates this silence, and what is it inferring from the silence it created?}

This distinction sharpens position one of the paper. Compliance theater is partly the failure to represent structural absences. But compliance theater's persistence is more substantially about produced absences: architectures that suppress challenge, then read the absence of challenge as evidence the architecture is correct. The dishonest middle is not merely a failure of representation. It is an active production of silence, interpreted as endorsement.

\subsection{The empty chairs in BSA/AML monitoring}
\label{sec:frame-as-method:bsa-aml-empty-chairs}

Six absent parties have interests in how a community bank monitors transactions for suspicious activity \parencite{ffiecBSAAMLManual}. Each is structural, produced, or both. We mark which.

\emph{The next victim of fraud whose pattern this transaction is part of.} Structural absence. When a transaction is part of an emerging fraud pattern, the institution's monitoring decision affects not the customer making the transaction but the customer who will be victimized in the next iteration. The next victim is fully absent from the transaction record and from the institution's customer relationship by the structure of time and the limits of prediction. They are represented in the monitoring architecture or they are not represented at all.

\emph{The customer whose account is frozen on a false positive.} Mixed: structural and produced. The customer is present in the transaction record but absent from the decision context that determined how their pattern was scored; that absence is structural, given that the model evaluates the transaction without consulting the customer. But the customer's silence after the freeze is partly produced: appeal processes are typically opaque, expensive, and slow, and many false-positive customers absorb the cost rather than contest it. The institution's apparent flagging accuracy is sustained in part by appeal-cost dynamics that suppress challenge from the falsely flagged.

\emph{The marginalized customer whose transaction patterns are flagged at elevated rates.} Mixed: structural and produced. Models trained on historical SAR data are at risk of inheriting historical flagging biases, especially where enforcement and reporting have differed across communities even when underlying risk has not \parencite{treasuryDeRisking2023}. Communities whose transaction patterns differ from the model's training distribution (recent immigrants using money-transfer services, cash-economy participants, members of religious or cultural communities with non-standard remittance patterns) can experience flagging rates that do not track underlying risk consistently across populations. The training-distribution gap is structural absence: the model was not trained to represent these populations as separate. But the appeal-cost dynamic is produced absence: the population most likely to be falsely flagged is also the population least equipped to contest the flag, and the institution's stable false-positive rate across populations is partly an artifact of differential appeal costs rather than evidence of unbiased flagging.

\emph{The institution itself, in the form of its future enforcement exposure.} Structural absence. The bank that fails to file required SARs faces regulatory action; the bank that over-files faces examiner skepticism about its risk calibration. Both are forms of enforcement exposure that compound over time. The institution's future self --- the bank in five years, when current monitoring decisions become the subject of an examination --- is an empty chair the present institution can either represent through architectural attestation or fail to represent through reliance on reconstructed reasoning.

\emph{The financial system, in the form of monitoring quality across institutions.} Structural absence. BSA/AML is a coordinated regulatory regime; one institution's monitoring inadequacy creates externalities for institutions whose monitoring is more rigorous, because illicit activity predictably gravitates toward weaker control environments. The system as a whole is an empty chair represented by inter-agency information sharing, FinCEN reporting, and cross-institution examination practice. The community bank's monitoring architecture either contributes to this representation or free-rides on it.

\emph{The criminal whose laundering depends on monitoring inadequacy.} Structural absence by intentional design. We name this empty chair explicitly because its absence in the architecture is intentional and structural, and because conflating ``we should not represent the criminal'' with ``we should not consider the criminal's interest in our threat model'' is a category error. The criminal's interest is the threat the architecture is designed against. The frame requires naming the chair we are designing to exclude. We also note the boundary case: the criminal and the marginalized customer can be hard to distinguish from behavior patterns alone, and architectures that conflate the two by design produce a particularly vicious form of produced absence: the marginalized customer flagged as criminal, whose appeal is expensive and whose silence after flagging is read as confirmation of suspicion.

The 230 control objectives do not generally describe controls in these terms. They specify governance outcomes for AI systems --- inventory attributes, oversight mechanisms, explanation practices, monitoring processes --- without naming every party whose interest an implementation might protect, and without distinguishing structural from produced absence. The empty-chair frame's first contribution is a vocabulary for examining those implementation choices: which chairs a control objective can represent, which chairs a particular implementation actually represents, and whether any remaining absence is given by the setting or produced by the implementation itself.

\subsection{Walking three examples through the frame}
\label{sec:frame-as-method:walking-three-controls}

The examples below apply four cross-sector FS~AI~RMF objectives to a community bank's BSA/AML transaction-monitoring system. They are not BSA/AML-specific requirements. For each example, we distinguish the objective's text from its implementation guidance and example evidence, and both from the bank implementation we ask the reader to consider. This distinction matters: an objective can support several implementations, and an implementation may support inferences the objective neither requires nor warrants. The objective identifiers and descriptions below are from the full Control Objective Reference Guide, version~1.0 \parencite{fssccFSAIRMF2026}.

\emph{Inventory: GV-1.6.1, AI Inventory Management.} The objective requires a regularly updated inventory of AI systems that includes purpose, data sources, responsible personnel, risk ratings, criticality, dependencies, compliance requirements, and decommissioning attributes. Its implementation guidance recommends a standardized template, assigned maintenance responsibility, periodic or trigger-based review, and access controls. Its examples of effective evidence include inventory records, ownership details, change histories, and audit trails.

Applied to transaction monitoring, a current inventory entry establishes something important: the institution has identified the system, assigned responsibility, and recorded specified lifecycle and risk attributes. It allows a future compliance officer or examiner to locate the system in the institution's governance landscape. It does not, by itself, establish that the listed data sources represent affected populations adequately, that the documented purpose matches operational use, or that a particular alert was well founded. Those are different propositions requiring different evidence.

The empty-chair question reveals the difference. The inventory directly represents the institution's future self and the future compliance officer by preserving organizational location and responsibility. It represents a customer affected by an alert only indirectly, through attributes that make later inquiry possible. A bank that treats the existence of a complete inventory entry as evidence that individual outputs are justified has asked the inventory to establish more than it can. Closing that gap requires evidence linked to the deployed system and the decision under review: the operative model version, relevant input lineage, validation results, and the record of how the output entered the decision process. The inventory is the index into that evidence, not a substitute for it.

\emph{Human oversight: GV-3.2.2, Monitoring Human-AI Oversight.} The objective requires mechanisms such as audits, performance evaluations, and feedback loops to assess human--AI interactions, including automation complacency and inappropriate overreliance. Its guidance asks institutions to examine how well operators detect issues, respond to alerts, and understand system limitations. Its examples include interaction audits, operator evaluations, feedback mechanisms, training, and ongoing assessment of oversight quality.

In a transaction-monitoring implementation, completed audits and operator evaluations establish that the institution examined specified aspects of human oversight. Measures of response quality, override patterns, or detected errors can show whether reviewers are engaging with alerts rather than merely accepting them. But a signed review record establishes only that a reviewer performed the recorded act under the institution's procedure. It does not, without further evidence, establish that the alert was substantively correct or that the reviewer independently evaluated every consideration material to the disposition.

This distinction serves both the falsely flagged customer and the next potential victim. For the former, a review signature is not evidence that the alert's basis was sound. For the latter, an override is not evidence that the reviewer noticed every risk the system surfaced. An implementation can narrow the gap by preserving what information was presented, what sources the reviewer consulted, whether the reviewer accepted or changed the recommended disposition, and which uncertainty or evidence triggered that action. Even this contemporaneous record has limits: it documents observable review activity, not unmediated access to the reviewer's mental process. The objective asks whether human oversight is effective; the implementation must specify which evidence makes that effectiveness assessable.

\emph{Explanation and decision context: MS-2.9.1, Model Explainability and Validation, and MS-2.9.2, Interpretability and Decision Context.} MS-2.9.1 calls for structured model explanations covering matters such as design choices, training-data sources, feature importance, and decision pathways, together with validation for stability, fairness, performance, bias, and limitations. Its example controls expressly include interpretable models and post-hoc tools such as SHAP and LIME, with documentation of methods, parameters, assumptions, and limitations. MS-2.9.2 requires outputs and decision rationales to be interpreted for their use context and intended audience. Its guidance and examples emphasize clarity, accessibility, human review, feedback, and testing the usefulness of explanations.

These objectives recognize several legitimate evidentiary functions. Model documentation can establish declared design choices and training sources. Validation can establish measured performance under specified tests. A SHAP or LIME report can characterize aspects of model behavior under the method's construction. User testing can establish whether an intended audience finds an explanation understandable. None of these results automatically establishes the others. In particular, clarity to a reviewer is not evidence that an explanation faithfully represents every computational dependency material to an output, and a technically faithful attribution is not evidence that the attributed model behavior is lawful or policy-conforming.

The empty-chair reading therefore asks which proposition matters to the affected party. A customer challenging a false alert may need to know which evidence supported the institution's action and whether that evidence was applied consistently with policy. An examiner may need to know the model's measured behavior across populations and whether the explanation procedure is stable under relevant conditions. A reviewer may need an account calibrated to the decision they are authorized to make. Treating one explanation artifact as satisfying all three needs produces a gap precisely where the objectives invite contextual implementation. Closing or disclosing that gap requires the institution to name the artifact, the proposition it supports, the method and assumptions that produced it, and the additional evidence needed for any broader inference.

\subsection{What the frame produced}
\label{sec:frame-as-method:what-the-frame-produced}

Four moves, in this section, the control objectives alone did not produce. First, an explicit accounting of which absent parties have interests in BSA/AML monitoring, including the chair we are designing to exclude. Second, a distinction between structural and produced absence, with different diagnostic moves for each. Third, a vocabulary for separating objective text, implementation evidence, and the further inferences drawn from that evidence. Fourth, a diagnostic question pair --- \emph{which empty chair does this implementation represent, and what silence does it produce that it then reads as endorsement?} --- that supplements the compliance question --- \emph{does this implementation satisfy the objective?}

These are not solutions to the problems the objectives address. They are a frame for interpreting the objectives and testing what their implementations establish. The frame's claim to organizing-principle status rests on whether moves of this kind generalize across controls and regulated decision contexts, which the remainder of this paper develops. What this section establishes is narrower: the frame makes visible a distinction among the required governance outcome, the artifact offered as evidence, and the conclusion drawn from that artifact. The structural-versus-produced distinction then asks why any missing evidence is absent and whose interests bear the consequence.

% Section 3: Verification Regimes as Taxonomic Move
%
% Status: first-pass integrated draft. Lending primary, underwriting as turn
% near the end. Section 2 took the unusual example (BSA/AML); this section
% takes the canonical one (lending) for load-management reasons articulated in
% planning document. Will be reworked.
%
% Citations needed: Rudin 2019, Bordt et al. 2022, Lanham et al. 2023, Turpin
% et al. 2023, ECOA / Reg B adverse action notice requirements, FCRA notice
% provisions, recent CFPB guidance on adverse action notices for AI-driven
% decisions, fair lending literature on disparate impact in pricing, recent FTC
% and CFPB statements on algorithmic pricing.

\section{Verification Regimes as Taxonomic Move}
\label{sec:verification-regimes}

Section~\ref{sec:frame-as-method} demonstrated that the empty-chair frame separates a governance objective, the artifact offered as evidence, and the conclusion drawn from that artifact. This section applies that distinction to consumer lending. AI applications in regulated decision contexts differ structurally according to the claims their verification regimes can support: which evidentiary objects the architecture exposes, how those objects bind to the decision, and what an affected party or examiner can test.

The worked example is the adverse action notice regime under the Equal Credit Opportunity Act and Regulation~B. The notice communicates specific reasons for a denial to an applicant who was absent from the decision but is legally entitled to an account of it. That function creates a precise verification question: what does the notice establish about the basis of the creditor's action, and what evidence allows the applicant or an examiner to test that relationship?

\subsection{The structural line}
\label{sec:verification-regimes:structural-line}

Recall the opening's framing: a verification regime lets a party test a material claim about a decision. Seven evidentiary objects recur in AI-supported workflows. Decision records preserve \emph{inputs and context}. SHAP can characterize \emph{functional feature dependence} under its value-function and background construction. LIME offers a \emph{local behavioral approximation} under its perturbation and surrogate choices \parencite{lundberg2017,ribeiro2016}. Attention or activation analyses expose selected \emph{internal computational states} \parencite{jain2019}. Chain-of-thought and other natural-language accounts are \emph{generated rationales}. Contemporaneous workflow records preserve observable parts of a \emph{human deliberative record}. Signatures and approval logs establish \emph{institutional authorization and review}. A system may expose several of these objects, but none acquires the evidentiary function of another merely because both are called explanations.

The methods therefore require method-specific readings. SHAP attributes an output under a specified construction; it does not by itself recover a human deliberative record or prove a complete causal account of model-internal computation. LIME approximates behavior locally, with a warrant bounded by its neighborhood and surrogate. An attention map exposes selected internal quantities, but attention weight alone is not a complete account of computational causation. Generated chain-of-thought is text whose faithfulness must be tested rather than assumed \parencite{lanham2023,turpin2023}. An ante-hoc interpretable model can expose its decision function more directly, while leaving separate questions about data provenance, policy validity, and deployment. A contemporaneous human record can show what was presented, consulted, recorded, and decided, but not every cognitive process. An institutional attestation proves that an authorized actor made a recorded commitment; whether the attested content is true depends on its evidence.

This taxonomy yields a structural line, but not a division between ``explanations'' and ``no explanations.'' An artifact cannot, by itself, verify a proposition about a different evidentiary object. Additional detail of the same kind does not close that gap. The gap closes when the architecture supplies an inspectable relation among the decision, the relevant object, the method that produced the artifact, and the proposition being tested. Section~\ref{sec:primitives} develops governance requirements for preserving and inspecting those relations. It does not promise transparent access to an anthropomorphic ``model reasoning'' object.

\subsection{Adverse action notices as the canonical case}
\label{sec:verification-regimes:adverse-action-notices}

The Equal Credit Opportunity Act requires creditors to provide adverse action notices to applicants who are denied credit, stating the specific reasons for the denial \parencite{regB1002_9}. The Fair Credit Reporting Act extends this to credit-based denials more broadly \parencite{fcraAdverseAction}. The regulatory intent is straightforward: a customer denied credit is entitled to know why, both to enable challenge of the denial and to inform the customer's understanding of their own creditworthiness. The empty chair the regime addresses is the denied applicant.

For a rule-based or low-dimensional scorecard, an institution may be able to trace a denial directly to the operative rule, threshold, input, and policy provision. Complexity changes the evidence problem, but it does not make every adverse action notice derived from a complex model unfaithful. Institutions may use feature attribution, local approximation, counterfactual analysis, model-specific decision paths, or combinations of methods. The CFPB does not permit complexity to excuse inaccurate or non-specific reasons \parencite{cfpbCircular2022_03}; the creditor must determine which method supports the reasons it communicates.

The notice itself establishes that the creditor communicated stated reasons. A method record can establish how those reasons were generated. Validation evidence can establish how the method behaved under specified tests. The further proposition --- that the communicated reasons accurately identify the factors material to this creditor's action on this application --- depends on the relation among those artifacts. Understandability is relevant but not sufficient; method fidelity is relevant but not sufficient; a plausible list is not the same as a tested binding to the action.

The empty-chair reading makes the evidentiary sequence visible. The applicant receives a legally required artifact. Their interest is represented substantively when the creditor can show why the artifact supports the proposition the law requires and when the applicant or examiner has a tractable way to challenge that showing. The frame does not declare every post-hoc method inadequate for every purpose. It asks the institution to name the purpose, the artifact, its warrant, and its limits.

The produced-absence dimension, picking up the distinction from Section~\ref{sec:frame-as-method}: the applicant who would have challenged the denial if challenge were tractable does not challenge. Adverse action notices are written in regulatory boilerplate, the reasons are typically generic (``credit history,'' ``debt-to-income ratio,'' ``insufficient income for amount requested''), and the applicant has no mechanism to interrogate whether these reasons accurately describe the model's behavior on their application \parencite{cfpbCircular2022_03}. The architecture produces silence (applicants who do not appeal because appeal is structurally unavailable) and reads the silence as evidence that the notice was adequate. The institution's position that adverse action notices satisfy ECOA / Reg~B is sustained partly by the fact that customers cannot effectively contest the notices' faithfulness, not by the notices being faithful.

\subsection{The taxonomic move}
\label{sec:verification-regimes:taxonomic-move}

The taxonomy divides applications by the verification claim their architecture can support.

\emph{AI that surfaces material for an independently accountable disposition.} Pattern recognition can surface applications for investigation; anomaly detection can flag unusual facts; consistency monitoring can identify departures from institutional norms. If an authorized reviewer receives the information and authority needed to make the disposition, the institution can preserve the decision inputs and context, the reviewer's observable actions, the recorded basis for the disposition, and the authorization. This does not make the human record infallible or expose every cognitive process. It produces an accountable institutional decision whose evidentiary components can be inspected.

\emph{AI that supplies the disposition and an account of it.} An underwriting model may generate a recommendation or decision, an explanation method may generate principal reasons, and a human may approve the result. This architecture is not unverifiable merely because it uses AI or a post-hoc method. Its burden is different: the institution must validate the method appropriate to its target, bind the resulting artifact to the operative model and case, define what the human review establishes, and preserve evidence for the broader substantive claim. A signature cannot supply a missing method-to-decision binding; nor does the presence of an explanation make such a binding impossible.

These are architecturally different, not simply points on a risk scale. The first locates the disposition and its recorded basis in an authorized human process supported by AI. The second locates material parts of the disposition in a model-supported process and therefore requires evidence across model behavior, explanation method, case context, policy, and review. ``Authority to override'' answers who may intervene. It does not answer what the reviewer could test, whether independent judgment was exercised, or what the approval establishes.

\subsection{What the categoricals foreclose, and what they do not}
\label{sec:verification-regimes:foreclosures}

The structural line forecloses specific \emph{claims}, not entire model classes. The following claims are not established without additional binding evidence:

\begin{itemize}
\item That an adverse action notice is faithful to the creditor's action solely because a post-hoc method generated plausible principal reasons and a human approved the notice.

\item That a decision-recommendation workflow received substantive human verification solely because an authorized reviewer signed or had formal override power.

\item That a retrospective justification reconstructs the decision-time information and deliberative record solely because no contradictory contemporaneous record exists.
\end{itemize}

The structural line permits stronger claims where the corresponding evidence is present, including:

\begin{itemize}
\item Ante-hoc interpretable models whose operative decision function, inputs, thresholds, and policy relationship are available for inspection, with separate validation of provenance and deployment.

\item Decision-support tools where AI surfaces material, an authorized human makes the disposition, and the architecture preserves the information presented, observable review activity, recorded basis, and authorization.

\item Model-supported decision architectures that validate an explanation method for its stated target, bind its output to the operative model and case, document its assumptions and limits, and preserve the other evidence needed for the substantive claim under review.
\end{itemize}

\subsection{A turn: pricing as produced-absence canonical case}
\label{sec:verification-regimes:pricing-turn}

The lending example demonstrates the categoricals at their canonical regulatory location, where the dishonest middle takes the form of an artifact (the adverse action notice) that misrepresents the reasoning behind a denial. This is not the only form the dishonest middle takes. We turn briefly to risk-based pricing in consumer lending (the determination of the interest rate, fees, and terms offered to an \emph{approved} applicant) to demonstrate that the frame produces interpretive leverage in regulatory spaces where the empty chair is not addressed by an artifact at all.

The Truth in Lending Act requires disclosure of a rate's mathematical structure (APR, finance charges, payment schedule); it does not require disclosure of why this particular rate was offered to this particular applicant \parencite{regZ1026}. Unlike the denied applicant, the priced applicant has no regulatory entitlement to a substantive explanation of their rate. The applicant receives a rate; the applicant accepts or declines; the institution proceeds.

The dishonest middle in pricing is more diffuse than in denial because there is no regulatory artifact equivalent to the adverse action notice. The pricing decision is not communicated to the customer as a reasoned determination; it is communicated as a take-it-or-leave-it offer. The customer accepts because they need credit and have no basis for negotiation. The institution reads the customer's acceptance as evidence that the rate was acceptable~\parencite{bartlett2022consumer}.

This is produced absence at architectural scale. The customer who would have challenged the rate if challenge were a thing they expected to receive does not challenge, because the architecture of risk-based pricing does not contemplate challenge as a normal part of the transaction. The institution reads the customer's silence (accepted offers, completed transactions, regulatory examinations finding no fair-lending issues) as evidence that pricing is fair. The architecture has produced the silence; the institution interprets the silence as endorsement.

The frame's diagnostic move on this configuration is the produced-absence move from Section~\ref{sec:frame-as-method}: \emph{what is the architecture doing that creates this silence, and what is it inferring from the silence it created?}

The architecture is producing pricing decisions through models whose reasoning is not communicated to the customer; structuring the transaction such that the customer cannot effectively negotiate without comparable offers they cannot tractably obtain; and producing aggregate pricing data that institutions analyze for fair-lending compliance, while the data itself is censored: customers who would have asked for justification if justification were available do not exist as a population the institution measures.

The institution then infers from this silence that pricing is fair. The customers who accepted are taken as evidence the rates were acceptable. The customers who declined are taken as having self-selected out. The customers who never applied because they expected unfavorable rates are not measured at all. The architecture has produced a population of accepted-and-silent customers and reads their silence as evidence of pricing fairness.

A common rejoinder: pricing models are statistically validated against outcomes (default rates by pricing tier, fair-lending disparate-impact analysis, ongoing monitoring \parencite{sr11_7,sr26_2}). This is true and important, and it is not a verification regime for the empty chair we are discussing. Statistical validation against outcomes verifies that the model performs as intended in aggregate; it does not verify that any individual customer's rate reflects their creditworthiness rather than model artifact. The aggregate empty chair (the regulator examining fair-lending compliance) is partly represented; the individual empty chair (the priced applicant) is not. The frame's contribution here is the aggregate-versus-individual distinction the regulatory regime currently elides.

\subsection{The taxonomic move, summarized}
\label{sec:verification-regimes:summary}

The lending example demonstrates the categoricals at the regulatory location where the dishonest middle is most directly addressable: where an artifact exists, where the regulation requires the artifact to represent the empty chair, and where post-hoc explanation methods have produced an artifact that does not. The pricing turn demonstrates that the frame extends to regulatory locations where the empty chair has been produced-absent by the regulatory regime itself, and that the frame proposes architectural primitives that would represent the chair the regulation does not currently require representing.

The remaining sections develop the diagnostic techniques that operationalize this taxonomic move (Section~\ref{sec:diagnostics}) and the categories of architectural primitive that empty-chair representation requires (Section~\ref{sec:primitives}). The frame's claim is that an honest AI governance framework must hold the structural line, both where the regulation requires it (as in adverse action notices) and where the regulation has not yet recognized that the line exists (as in risk-based pricing). Whether this claim is sustainable in practice is the subject of Section~\ref{sec:uncertainties}'s honest uncertainties.

% Section 4: Diagnostic Techniques

\section{Diagnostic Techniques}
\label{sec:diagnostics}

The empty-chair frame produces interpretive structure (Section~\ref{sec:frame-as-method}) and a taxonomy of evidentiary claims (Section~\ref{sec:verification-regimes}). Three gaps recur across those sections: documentation may be treated as proof of the activity it describes; an explanation artifact may be treated as proof of a different evidentiary object; and human approval may be treated as substantive verification of material the reviewer could not test. The artifacts themselves can be genuine and useful. The failure occurs in the unsupported inference from the exposed artifact to the unexposed proposition. We name the recurring pattern \emph{silence-manufacture}. The three diagnostic techniques that follow are designed to expose the missing relation and the absent parties who depend on it.

We frame the techniques from the architect's perspective because the techniques are most powerful when they constrain construction rather than evaluate existing artifacts. An architect who applies them builds systems that are diagnosable; an examiner who applies them to an undiagnosable system can only document that diagnosis is unavailable~\parencite{brooks1974mythical,wing2008computational}. The architectural moment is where empty-chair representation is structurally determined, and the diagnostic techniques are most valuable as design-time discipline. The examination-time application of these same techniques falls out as a consequence and is treated below as such.

\subsection{Silence-manufacture as the underlying structural pattern}
\label{sec:diagnostics:silence-manufacture}

\subsubsection{Definition}
\label{sec:diagnostics:silence-manufacture:definition}

We define \textbf{silence-manufacture} as a structural pattern with four parts: access to an evidentiary object material to a decision is unavailable or suppressed; a different artifact is produced; the artifact is treated as establishing the unavailable object; and the absence of contradictory evidence is taken as support for that inference. A metric, document, signature, or explanation becomes a substitute only when it is asked to warrant a proposition beyond the object it exposes.

Silence is the \emph{enabling condition}; substitution is the \emph{active mechanism}. The artifact is real. A SHAP value is a real numerical computation. An audit document is a real institutional production. A signed adverse action notice is a real legal instrument. Each can support claims appropriate to it. Silence-manufacture begins when the institution treats the artifact as evidence of a different proposition without making the relation inspectable, while the architecture leaves affected parties unable to produce the evidence that would test that relation. The concept therefore criticizes neither documentation nor explanation as such. It identifies a specific inference gap.

\subsubsection{Relationship to Goodhart's law}
\label{sec:diagnostics:silence-manufacture:goodhart}

Silence-manufacture is the \emph{dual} of Goodhart's law \parencite{goodhart1984problems}, not strictly its inverse. The relationship is rhetorically useful for orienting readers (Goodhart provides a known landmark) but the load-bearing name is silence-manufacture, not extended-Goodhart.

Goodhart's law describes optimization pressure on a metric corrupting the metric: when a measure becomes a target, it ceases to be a good measure. The pressure lands on the visible; the visible breaks. Silence-manufacture describes suppression of evidence needed to test the inference drawn from a visible artifact: silence is treated as confirmation even though the architecture helps produce it. The pressure lands on what cannot readily be observed; the artifact appears adequate because the evidence needed to contest its broader use is missing.

The two patterns are related through structural symmetry rather than direction-reversal. Goodhart describes corruption of metrics; silence-manufacture describes an evidentiary substitution that can leave the metric intact. A regulator may improve a visible measure while leaving unchanged the architecture that prevents an affected party from testing what the measure is taken to establish. We return to this possibility below as a proposition generated by the frame, not as a universal law of reform.

\subsubsection{Relationship to decoupling and the audit society}
\label{sec:diagnostics:silence-manufacture:decoupling}

Silence-manufacture has close neighbors in organizational sociology. Meyer and
Rowan's account of decoupling describes formal structures maintained apart from
operational activity; Power's audit society and Strathern's account of audit
culture describe verification practices that can displace the substance they
were meant to secure \parencite{meyer1977,power1997,strathern1997}. The broad
phenomenon of an artifact standing in for activity is therefore not new, and we
do not claim it is.

The narrower contribution here is a diagnostic sequence. What evidentiary
object does a claim require? What feature of law, architecture, interface,
method, data, or cost makes that object unavailable? What artifact is offered
instead, what inference is drawn from it, and which absent party depends on the
inference? This sequence does not presume organizational bad faith or even a
coupled alternative that was feasible. It can therefore distinguish ceremonial
decoupling from cases in which the relevant silence is legally required,
architecturally produced, or methodologically unavoidable.

\subsubsection{Mapping the three earned positions onto silence-manufacture}
\label{sec:diagnostics:silence-manufacture:mapping-positions}

The three positions expose the same evidentiary sequence at different levels,
but they do not depend on identical technical mechanisms. Position~1
(governance versus governance theater) is the enclosing institutional frame;
Positions~2 and~3 identify particular explanation and accountability failures
within it. Their common structure is the inference from an exposed artifact to
a proposition that the artifact, without an inspectable binding, does not
establish.

\emph{Position 2: explanation artifacts do not verify evidentiary objects they
do not expose.} The relevant limitation depends on the method. SHAP attributes
a prediction under a specified feature-coalition formulation; LIME fits a local
surrogate around an instance \parencite{lundberg2017,ribeiro2016}. Neither is,
without an additional warrant, a transcript of internal computation or a human
decision process. Attention weights have likewise been shown not necessarily to
support the causal interpretations assigned to them \parencite{jain2019}.
Chain-of-thought studies provide a different kind of evidence: generated traces
can omit influential cues or vary independently of answers in studied settings
\parencite{lanham2023,turpin2023}. Slack et al. show that post-hoc methods can
also be manipulated to conceal biased classifier behavior \parencite{slack2020}.
These results do not make the methods interchangeable or useless. They show why
the proposition asserted must be matched to the artifact actually produced.

\emph{Position 3: review of an explanation is not verification of an unexposed
decision process.} A signature can establish authorization, completion, and the
identity of a reviewer. It supports a stronger claim only when the reviewer had
access to evidence capable of testing that claim and the record shows what was
tested. Research on sycophancy and evaluator-facing optimization illustrates
why approval signals can diverge from the qualities they are taken to measure
\parencite{sharma2023,perez2023}; it does not establish that every approval does
so. The regulator's question is therefore about reviewer access, competence,
and recorded basis, not the mere presence or absence of a signature.

\emph{Position 1: governance documentation is evidence of documented
commitments and activities, not automatically of operational effect.} Policies,
inventories, dashboards, and audit records can be indispensable governance
artifacts. The inference gap appears when their existence is treated as proof
that deployed behavior conforms to them, although the relevant operations were
not independently observed or tested. Here again, the remedy is not less
documentation but a declared relation between the artifact, the proposition,
and the evidence capable of defeating the proposition.

The unification is diagnostic rather than causal. Across the three levels, an
examiner can identify the proposition at issue, the available artifact, the
claimed binding, and the evidence that could disconfirm it. That common sequence
is useful even when the causes of unavailability and the appropriate remedies
differ.

\subsubsection{What silence-manufacture predicts}
\label{sec:diagnostics:silence-manufacture:predictions}

The frame suggests three propositions that distinguish it from a generic
Goodhart-shaped critique. We state them with observations that would count
against them.

\emph{Proposition 1: metric-layer reform can leave an inference gap intact.}
Tighter thresholds, controls, or documentation requirements can improve a
visible artifact without making independently testable the broader proposition
for which it is used. The proposition would be weakened where metric-only
reforms reliably make that broader proposition independently testable, without
changes to evidence capture, access, or binding. Section~\ref{sec:primitives}
develops architectural responses; Section~\ref{sec:uncertainties} marks the
political-economy assumptions on which adoption depends.

\emph{Proposition 2: inquiry can partition around unavailable evidence.}
Observers with different access may form durable, internally coherent accounts
when evidence capable of resolving their disagreement is unavailable. The
proposition would be weakened if the disagreement persisted after the relevant
evidence became mutually available and its provenance could be tested. More
research is not necessarily curative; the diagnostic question is whether the
researchers are able to observe the object over which they disagree.

\emph{Proposition 3: public narratives can inherit the shape of unavailable
evidence.} The BSA/AML prohibition on disclosing suspicious-activity reports
creates a concrete setting in which customers and outside observers may lack
access to an institution's account of a closure \parencite{usc31_5318g2}.
Anthony's finding that 35 of 8,361 account-closure complaints mentioned
political or religious motivations illustrates the need to distinguish a
salient narrative from the observed complaint distribution
\parencite{anthony2026}; it does not prove what caused every closure or every
narrative. The proposition would be weakened if access to the previously
unavailable evidence did not materially change the explanations formed by
affected parties or public observers.

These are research propositions, not results established by this position
paper. Their value is that they identify observations capable of refining or
defeating the frame at three sites: regulatory reform, analytical inquiry, and
public discourse.

\subsection{Three diagnostic techniques}
\label{sec:diagnostics:three-techniques}

The diagnostic techniques developed below are silence-breaking instruments. Each technique is paired to a class of silence the architecture produces, and the technique's value is in what it makes audible. We claim three techniques because they emerge cleanly from the frame's structural commitments and produce design constraints directly applicable to AI governance system construction. Other techniques are possible and likely; these three are the ones the frame produces first, and together they cover the silence-types the three earned positions identify.

\subsubsection{Technique 1: Empty-chair enumeration as design pre-condition}
\label{sec:diagnostics:enumeration}

Before specifying any architectural primitive, the architect asks: \emph{which empty chairs is this system supposed to represent, and what does representation require for each?} The question precedes any technical decision. An architecture whose empty-chair attribution is left implicit will produce primitives that serve some chairs and fail to serve others, and the failure will not be visible in the architecture's documentation because the documentation will describe what the system does, not what it leaves out. The silence the technique breaks is the silence around \emph{which absences the architecture is making consequential}: an architecture that has not enumerated its empty chairs is suppressing the question of which chairs it serves, not answering it.

Section~\ref{sec:frame-as-method} demonstrated this technique applied analytically to BSA/AML transaction monitoring. The same move applies to design. The architect designing an AI governance system for a community bank's lending operation asks first: \emph{whose interest is this system structurally protecting?} The answer is rarely on the page initially; it has to be derived from the regulatory regime, the institution's risk profile, and the categories of decision the system will affect. Making the derivation explicit at design time is the technique's primary work.

The technique produces three outputs at design time:

First, an enumeration of empty chairs the system affects. This is the analytical move from Section~\ref{sec:frame-as-method}, applied at the moment of architectural commitment.

Second, a mapping from empty chairs to architectural requirements that representation imposes. Different empty chairs require different architectural support. The future compliance officer needs durable, queryable institutional knowledge. The denied applicant needs a testable account of the action's basis. The marginalized customer needs disparate-impact monitoring with population decomposition. The institution's future self needs tamper-evident temporal capture. Each chair generates a set of requirements; the requirements may conflict, and the conflicts are themselves design information.

Third, a constraint on architectural choices: any primitive that fails to satisfy at least one chair's requirement and creates costs without serving any chair is architectural overhead. Any primitive that satisfies one chair's requirement at the cost of failing another's is a design tradeoff that requires explicit attention. The architecture is not a set of capabilities; it is a set of empty-chair representations, with capabilities chosen as the means.

The technique shifts architecture from feature selection to representation. Without it, systems can accumulate explainability features, audit trails, and monitoring dashboards without identifying which empty chairs those capabilities serve. With it, architects start from the chairs and select capabilities as the means of representing them. The criterion makes those choices tractable; without it, they can reduce to keeping up with vendor offerings.

The examination question is concise: \emph{Which absent party depends on this artifact, and what evidence would show that the implementation represents that party's interest?}

\subsubsection{Technique 2: Designing for pages missing from the ledger}
\label{sec:diagnostics:pages-missing}

A well-designed attestation regime makes the institution's recorded commitments
and decision evidence inspectable. But attestation cannot be uniformly
comprehensive; cost and operational reality require choices about what to attest
and what to leave un-attested. The architect's diagnostic move at design time is
to make these choices explicitly and to design the attestation regime such that
\emph{what is left out is itself informative}. The silence the technique breaks
is the silence between the institution's posture and its documentary record:
under ambient absences, posture is unreadable; under documented absences,
posture becomes a pattern an examiner can read.

We call this technique \emph{designing for pages missing from the ledger}. The metaphor is deliberate: an accounting ledger with missing pages is more diagnostic than a forged ledger if the pattern of missing pages is itself interpretable. A ledger where missing pages are random reveals nothing about the institution's posture; a ledger where missing pages cluster around decisions favorable to the institution and unfavorable to challenged customers reveals the institution's posture clearly.

The technique works as follows. The architect determines, for each category of decision the system handles, what attestation will and will not be produced. The determination is documented. The institution's choice not to attest a category is itself attested, with the rationale captured. An examiner reading the system's attestation later can read both the present attestation (what was captured) and the meta-attestation (what was deliberately not captured, and why). The pattern of present and absent attestation is interpretable because the absences are themselves documented choices rather than oversights.

This is a design discipline, not a construction discipline. Most attestation systems are built bottom-up: a capability is implemented, attestation for that capability is added, and the absence of attestation in other areas is ambient. Designing for pages-missing inverts this: the architect specifies in advance the full domain of decisions, marks which will and will not be attested, and ensures the system's structure makes the absence visible rather than ambient.

The technique produces a design output: a manifest of attested and un-attested decision categories, with rationale, that travels with the system as part of its architectural documentation. An examiner reading the manifest can ask whether the un-attested categories are operationally explicable (resource constraints, low-risk decisions deprioritized) or whether they cluster in ways that reveal a posture (decisions favorable to the institution attested, decisions unfavorable to challenged customers not). The architect who has produced the manifest knows in advance which questions the manifest will answer and which it will not, and can adjust the attestation scope accordingly.

The examination question is: \emph{Which decision categories lack contemporaneous evidence, was that absence declared in advance, and what pattern do the absences form?}

\subsubsection{Technique 3: The produced-silence question pair as design constraint}
\label{sec:diagnostics:question-pair}

Section~\ref{sec:frame-as-method} introduced the distinction between structural and produced absence. The diagnostic technique on produced absence, applied at design time, is a question pair the architect asks of every interface, every documentation regime, every customer-facing artifact:

\emph{What is this design doing that creates silence?}

\emph{What is this design inferring from the silence it creates?}

This is the silence-manufacture-direct technique. The first two techniques operate on silences that the architecture has implicitly allowed to form (empty chairs not enumerated, decision categories ambient about their attestation status); this technique operates on silences the architecture \emph{actively produces}, then \emph{reads as endorsement}. The two-part question maps directly onto the two parts of silence-manufacture: the first question surfaces the architecture's silence-production, the second surfaces the substitute artifact and the inference being drawn from suppression.

A design that asks reviewers to write justification after deciding leaves the
decision-time information and observable deliberative record unavailable, then
infers from the absence of contradicting evidence in the post-hoc justification
that the documentation is faithful. The architect applying the question pair at
design time recognizes the silence-production and chooses differently,
capturing decision context as the reviewer engages with it rather than asking
for narrative reconstruction.

A design that produces adverse action notices in regulatory boilerplate creates silence around the customer's effective interrogation of the notice, then infers from the absence of customer challenges that the notice is adequate. The architect applying the question pair at design time recognizes the boilerplate's interrogation-resistance and chooses differently, producing notices the customer can effectively interrogate, even if doing so requires architectural work the regulation does not currently require.

A design that prices loans without communicating substantive justification to the customer creates silence around the customer's standing to negotiate, then infers from acceptance rates that pricing is fair. The architect applying the question pair at design time recognizes the absence of customer-side justification mechanisms as a design choice with consequences, and weighs whether the consequences are acceptable given the empty chair the design leaves unrepresented.

The question pair, applied at design time, has the same structure as at examination time but different effect. At examination time, it diagnoses what the architecture has already done. At design time, it constrains what the architecture will do. The architect who anticipates the diagnostic question pair builds a system that can answer it; the architect who does not builds a system that the question pair will diagnose negatively.

This technique is the most contestable of the three because it requires the architect to design for absences they are choosing to allow. An architect may legitimately respond that some silences are structural rather than produced, that customers who don't challenge are satisfied, that reviewers who don't surface contemporaneous reasoning don't have any to surface. The architect may be correct in some cases. The technique's value is in making the question explicit at design time so the architect's choice is conscious rather than ambient.

The examination question is: \emph{What access did the architecture remove, what artifact replaced it, and what inference is being drawn from the replacement?}

\subsection{The techniques as compounding design constraints}
\label{sec:diagnostics:compounding}

The three techniques compound at design time. Empty-chair enumeration breaks the silence around which absences the system is structurally responsible to. Designing-for-pages-missing breaks the silence between institutional posture and documentary record. The produced-silence question pair breaks the silence the architecture itself manufactures and reads as endorsement. In combination, the techniques do not merely produce a checklist of design requirements; they produce a design discipline whose effect is to replace the architecture's manufactured silences with inspectable evidence at the points where silence-manufacture would otherwise sustain an unchallenged substitute.

This is structurally different from feature-driven architecture. The system's capabilities are chosen as means of empty-chair representation; the system's attestation regime is designed for pattern-readability; the system's interfaces are scrutinized for silence-production at the moment of design rather than at the moment of examination. The architecture that results is one that an examiner can read; the architecture without these techniques is one that an examiner cannot read regardless of skill.

This is the operational claim of the section: the empty-chair frame, developed into design-time diagnostic techniques against silence-manufactured artifacts, produces architectural discipline that makes the resulting systems substantively examinable. The frame is not an examination methodology that compensates for inadequate architecture; it is an architectural discipline that produces examinable systems. The examination-time application of the techniques follows from the architectural use, not the other way around.

\subsection{Examination as consequence}
\label{sec:diagnostics:examination}

The same three techniques, applied by an examiner to an existing system, produce assessment leverage the FS~AI~RMF's checklist does not provide. The examiner asks: \emph{which empty chairs does this system claim to represent? What does its attestation pattern reveal about its actual representation? What silences does its architecture produce, and what does the architecture infer from those silences?} These are the same questions the architect asked at design time, redirected from ``what should I build'' to ``what was built.''

The asymmetry between design-time and examination-time application is important. A system designed with the techniques in mind answers all three questions; an examiner reading such a system can assess whether the answers are substantive or theatrical. A system designed without the techniques in mind may be unable to answer the questions at all; an examiner reading such a system can only document the unavailability of substantive assessment.

The implication: the techniques' regulatory effect depends on adoption at the architectural level, not just at the examination level. An examiner who applies the techniques to a system designed without them produces findings of inadequacy, but the inadequacy is structural: it cannot be remediated by additional examination, only by re-architecture. The frame's claim is that AI governance systems for community banking should be architected to be diagnosable, and the diagnostic techniques developed here should shape the architecture rather than only its evaluation. The development of examination-time methodology is the work of Paper~3; this paper's diagnostic register is architectural because the architectural moment is where the techniques have most leverage.

\subsection{Implications for architectural primitive selection}
\label{sec:diagnostics:primitive-implications}

Section~\ref{sec:primitives} develops the categories of architectural primitive that empty-chair representation requires. The diagnostic techniques developed here constrain that development: a primitive is valuable in proportion to its support of design-time diagnostic discipline against silence-manufacture. Tamper-evident temporal capture is valuable because designing-for-pages-missing requires the ability to attest both presence and absence of records, and because the silence-manufacture pattern at the institutional level is sustained by retroactive narrative substitution that contemporaneous capture forecloses. Three-dimensional confidence characterization is valuable because empty-chair enumeration requires distinguishing decisions where the architect provided genuine evidentiary support from decisions where the architect produced narrative without substantive grounding, which is precisely the distinction silence-manufacture exists to obscure. The primitive categories are not generic technical capabilities; they are capabilities selected because the design-time techniques require them, and the techniques require them because silence-manufacture is the failure mode they are constructed against.

This is the section's transition into Section~\ref{sec:primitives}. The techniques operationalize the frame at design time; the primitives operationalize the techniques. Each layer of the architecture is justified by what the layer above it requires of it, and the layer above the techniques is the structural pattern they are constructed to break.

% Section 5: Provisional Evidentiary Capabilities

\section{Provisional Evidentiary Capabilities}
\label{sec:primitives}

The diagnostics in Section~\ref{sec:diagnostics} generate five provisional
evidentiary capabilities. They are design hypotheses, not universal
requirements or a complete governance architecture. A capability may be
supplied by a vendor product, a bank-controlled system, a contractually required
export, an independent assessment, or a compensating process. The institution
remains responsible for claims it makes with vendor-supported evidence, but
responsibility does not imply unilateral power to alter the vendor's system.

Each capability is stated by the object it should expose and the limit on what
that object establishes. An examiner need not personally perform cryptographic
verification or model forensics: the review may inspect bank or third-party
validation, provenance, controls, and selected underlying evidence. Which
pathway is proportionate depends on the claim, risk, authority, and available
implementation.

\subsection{Capability 1: Revision-aware temporal capture}
\label{sec:primitives:temporal-capture}

The first capability preserves distinct, time-bound records rather than a
single mutable narrative. Its decision-time record captures the inputs and context
presented, observable actions, stated uncertainties, and provisional reasons.
Later records may add reflective interpretation or newly discovered reasons;
still later records may authorize, correct, or reassess the decision under
policy. Each layer carries its own timestamp, provenance, and binding to the
decision and prior record. Later reflection may contribute genuine
understanding, but it must not overwrite or impersonate the earlier account.

Tamper-evident commitments, controlled audit logs, signed records, or
independently maintained timestamps are possible implementations. A bank may
operate them or contract for them. An examiner can inspect validation of the
binding and selected records rather than reproduce the mechanism. Temporal
capture establishes when specified content was committed and whether revision
is detectable; it does not establish that the content was true, complete, or a
transparent account of human cognition.

\subsection{Capability 2: Evidence binding and exportability}
\label{sec:primitives:evidence-binding}

The second capability makes a declared relationship among a decision, artifact,
method, model and data versions, policy, review, and authorization traversable.
Inspection should be possible backward from a decision to its supporting
material and, where appropriate, forward from material to affected decisions.
The export must preserve identifiers, provenance, configurations, and relevant
validation rather than reduce the relationship to a product dashboard.

This is particularly important for community banks using vendor-controlled
systems. A bank must be able to obtain, directly or by contract, evidence
sufficient to support the claims it is obligated to make. A vendor may provide
a structured export; a bank may maintain external bindings; an independent
assessor may validate proprietary components; or a compensating control may
narrow the claim. Product access without inspectable evidence leaves the
coordination gap in place.

Binding an explanation to a decision establishes association and provenance.
It does not prove causal influence, method fidelity, or evidentiary sufficiency.
Those are separate propositions tested with method-appropriate evidence.

\subsection{Capability 3: Non-collapsing uncertainty representation}
\label{sec:primitives:confidence-characterization}

The third capability preserves materially different reasons for uncertainty
instead of forcing them into one probability or confidence label. Depending on
the decision, a useful representation may distinguish supporting evidence,
counterevidence, conflict among sources, model or measurement uncertainty,
missing information, and ambiguity about applicable policy. A single value can
be well calibrated for one target while remaining silent about which of these
conditions produced it.

The paper does not prescribe a fixed number of dimensions or one mathematical
formalism. A structured record, interval, set-valued output, qualitative state,
or combination may be appropriate. The functional test is whether material
support, counterevidence, conflict, and insufficiency survive reporting and can
be bound to the resulting action. Vendors may expose native uncertainty and
provenance; banks must decide how policy treats the represented state; and
validation should test whether the representation is stable, meaningful, and
used as declared. The representation does not itself establish that the
underlying evidence is accurate.

\subsection{Capability 4: Technical-change monitoring}
\label{sec:primitives:technical-change}

Technical-change monitoring asks whether the empirical system has changed.
Relevant objects include input distributions, relationships between inputs and
targets, measured performance, calibration, subgroup behavior, and the model or
data pipeline itself. Data-drift tests, concept-drift analysis, outcome
monitoring, version comparison, and targeted validation answer different
questions; no shared detector is presumed.

The likely implementation path combines vendor telemetry and change notices
with bank or third-party validation against the institution's operational
population. Examiners can inspect the monitoring design, thresholds, escalation
records, and selected results. A detected change directs investigation; it does
not establish cause, harm, or noncompliance.

\subsection{Capability 5: Institutional-conformance monitoring}
\label{sec:primitives:institutional-conformance}

Institutional-conformance monitoring asks a different question: whether
operational practice continues to match declared policy, authority, and control
design. Its evidence includes workflow and override records, access and
authorization histories, exceptions, complaints, sampled case files,
attestation-scope reconciliation, and changes in policy interpretation.
Statistical drift may be absent while practice diverges from policy, or present
while practice remains conformant.

The bank ordinarily owns the policy relationship, while vendors must expose the
events and configurations needed to test it. Internal audit, compliance,
independent review, affected-party evidence, and examiner-selected samples can
challenge a self-selected institutional account. A discrepancy establishes a
conformance question and may support a finding where applicable authority and
evidence do so; the monitoring capability alone does not determine legal
significance.

\subsection{What the capabilities contribute}
\label{sec:primitives:compound}

Together, the five capabilities can make the diagnostics answerable:
revision-aware capture preserves temporal layers; binding and exportability make
declared relations inspectable; non-collapsing representation preserves the
kind of uncertainty at issue; and the two monitoring capabilities distinguish
technical change from institutional divergence. Multiple coherent
implementations may exist, and a compensating process may sometimes provide the
required evidence more proportionately than new infrastructure.

Their presence does not establish that every affected interest is represented,
and their absence does not prove that representation is impossible. Cost,
security, confidentiality, records law, and vendor leverage constrain the
design space. The practical question is whether a proposed implementation makes
the relevant artifact-to-inference relation more inspectable at a proportionate
cost. That question remains empirical.

% Section 6: Honest Uncertainties

\section{Honest Uncertainties}
\label{sec:uncertainties}

The preceding sections bound their claims where they arise. This section does
not retract those claims. It collects the remaining uncertainties that matter
to a regulator, architect, or examiner deciding how to use the frame.

\subsection{Scope of the frame}

\textbf{Empty-chair enumeration and weighting remain open.} The six chairs in
the BSA/AML example are illustrative, not exhaustive. Practitioners may identify
different affected parties or assign different weight to their interests. The
frame contributes a discipline---name the absent party, the interest at stake,
and the evidence of representation---but it does not supply a complete roster
or a weighting rule. Disagreement is most useful when those choices are made
explicit and can themselves be examined.

\textbf{Structural and produced absence may not exhaust the possibilities.}
Some absences are mixed, and practice may reveal additional categories. The
distinction is useful where it changes the diagnostic response: a structural
absence may require representation or a proxy, while a produced absence calls
for scrutiny of the design that produced it and the inference drawn from it.
The paper does not claim that every case fits cleanly into one category.

\textbf{Generalization is limited to the contexts examined.} The paper develops
the frame through community-banking examples in BSA/AML monitoring, lending,
and pricing. Application to insurance, healthcare, criminal justice, or social
services is a research proposition, not a result. Even within finance, evidence
that a mechanism appears on one dataset or product does not establish transfer
to another. Sectoral use should begin with new empty-chair enumeration, legal
analysis, and validation of the relevant evidentiary relationships rather than
directly importing the examples here.

\subsection{Conditions of examination}

\textbf{Adversarial and cooperative settings differ.} Examination includes
both: an examiner may test a claim against an institution's incentives, and the
same parties may cooperate to resolve a shared safety or compliance problem.
The paper's strongest warnings about explanation and self-verification concern
settings where incentives are opposed or where a claim must survive challenge.
In cooperative settings, explanation artifacts can support discovery and
communication even when they would not independently verify the proposition at
issue. The method and claimed inference should therefore be evaluated in the
actual examination context.

\textbf{Application requires resources and coordination.} The three diagnostic
questions demand time, expertise, access, and institutional support. Not every
examiner can apply them to every system, and a sophisticated institution can
design formally compliant artifacts that remain difficult to interrogate.
Their practical effect may depend less on exhaustive case-by-case use than on
making the questions predictable enough that institutions design systems able
to answer them. This paper does not establish whether supervisory organizations
can provide the necessary capacity or coordinate a consistent practice.

\textbf{Affected-party access is not uniform.} The ontology gives affected
interests standing in the analysis; it does not imply that every affected party
can or should inspect every artifact. Suspicious-activity confidentiality,
personal data, cybersecurity controls, trade secrets, and the rights of other
customers can limit direct disclosure. Representation may therefore travel
through several contest pathways: a customer-facing explanation and correction
process, protected access for an examiner, independent testing of a proprietary
method, aggregate reporting, or an advocate with appropriate standing. These
pathways are not equivalent. Independent validation can test a bounded
proposition while leaving the affected party unable to identify an error in
their own case; a clear notice can enable challenge while revealing little about
population-level behavior. Applying the frame requires stating who can access
which object, what proposition that access permits them to test, and which
unresolved interest remains represented only indirectly. Whether a particular
access restriction is justified is a legal, security, and proportionality
question outside the ontology itself.

\textbf{Representation is not participation or remedy.} Indirect inspection by
an examiner or independent assessor may represent an affected interest in an
evidentiary inquiry while leaving the affected person without direct voice,
access, or relief. Existing complaint, appeal, correction, and review processes
may permit participation or contest, but this paper neither creates those
processes nor establishes when law should require them. Where none exists, the
ontology can make the absent pathway, the institution able to create it, and the
party bearing its absence legible. That diagnosis is not itself a remedy.

\subsection{Architectural boundaries}

\textbf{The five capabilities are provisional and non-exhaustive.}
Revision-aware temporal capture, evidence binding and exportability,
non-collapsing uncertainty representation, technical-change monitoring, and
institutional-conformance monitoring are generated by this frame's diagnostics.
They are neither a complete governance architecture nor proven necessary for
every deployment. Alternative technical, contractual, independent-assessment,
or compensating-process implementations may answer a diagnostic question more
simply. Implementation may also reveal additional needs.

\textbf{Operational tractability is unproven.} The capabilities may impose storage,
integration, validation, and review costs. Those costs may be material for a
community bank and may interact with confidentiality, retention, cybersecurity,
and vendor constraints. Banks remain responsible for their claims but may lack
leverage to obtain evidence from core-system and AI vendors. The paper has not
built or evaluated a complete implementation, compared alternatives, or shown
that benefits justify costs. The capabilities are design hypotheses with
explicit evidentiary functions, not an implementation mandate.

\textbf{Vendor dependence creates a two-level theory of effect.} A community bank may
identify an evidence gap correctly and still be unable to alter a core platform,
obtain method provenance, or negotiate a bespoke export. Large vendors can
spread implementation cost across institutions, but they also choose product
roadmaps and contractual defaults under incentives different from those of one
bank, examiner, or affected customer. One bank can make the dependency visible,
narrow its claim, evaluate available processes, and record the residual risk.
The collective pathway aggregates those limits: examiners identify unsupported
claims within current authority; regulators and trade groups observe recurrence;
and vendors receive a more consistent capability signal. This paper does not
establish that collective signals will overcome switching costs, concentration,
proprietary constraints, or fragmented regulatory expectations.

Nor does vendor implementation end the inquiry. A standardized export may omit
locally material context; independent validation may be too narrow or too
expensive; a compensating process may shift work to already constrained bank
staff. Confidentiality and cybersecurity can also justify limiting access that
would otherwise improve contestability. The relevant comparison is among
feasible evidence pathways and their residual gaps, not between perfect access
and presumed noncompliance.

\subsection{Who moves first}
\label{sec:uncertainties:first-move}

The largest uncertainty is not whether any participant would find the frame
useful. It is that each participant's first move appears to depend on another's.
A community bank converting silent gaps into written disclosures before an
examination has reason to want evidence that examiners will treat a disclosed
gap with a compensating process as adequate posture rather than as a
self-reported finding. A vendor asked to fund an evidence-export capability has
reason to want evidence that examiners are issuing requests in this shape, since
capability purchases follow examinations that institutions could not answer. An
examiner has authority to test a claimed inference but must be able to name the
statute, regulation, or supervisory expectation under which each evidentiary
object may be compelled, and to state what may be supportably concluded when an
institution answers that the object is vendor-proprietary and unavailable. A
regulator considering guidance must demonstrate workability for a small
institution, which requires evidence about burden that only institutional use
would generate.

Those four dependencies form a cycle. The paper cannot resolve it by asserting
that collective action will aggregate individual practice; that assertion is
what the cycle calls into question.

Two artifacts would break it, and neither requires any participant to act
first. The first is an authority crosswalk: for each evidentiary object this
paper names, the existing authority under which its production may be sought,
and the inference available when it is withheld as proprietary. The second is a
mapping from selected FS~AI~RMF control objectives to the records a conforming
institution would necessarily hold, which would let a party without discovery
request documents by name rather than by description. Both are desk work
against existing law, guidance, and the framework text. Neither is supplied
here. Their absence is the most consequential gap in this paper, and it is a
gap in this paper rather than in the coordination problem it describes.

A preliminary attempt at the first suggests the crosswalk will not divide the
objects into compellable and non-compellable. Supervisory authority is not
confined to inspecting records that happen to exist: reporting powers can
require statements and special reports prepared on demand, recordkeeping
regulation can mandate that records be generated prospectively, and an
examiner's own analysis of produced data is the examiner's work product rather
than something compelled from the institution. The line that does hold is
narrower and concerns time. An institution can be required to answer for a
missing record, to prepare a report, or to build a prospective control. It
cannot be required to produce, after the fact, an object that would exist only
if the decision process had preserved it contemporaneously.

That line cuts through the objects unevenly. Feature attributions, local
approximations, and internal-state observations are in principle recomputable
later. What is not recomputable is the binding: which model version and
configuration, under which policy, actually produced a particular decision on a
particular file. Where that binding was not captured at decision time, the
recomputable objects lose their referent, and the question of whether they
could have been compelled becomes moot --- there is nothing to compel them
against. Decision-time context and the record of what a human reviewer was
actually shown are subject to the same limit and have no recomputable
substitute at all.

The practical consequence is that one narrow record carries most of the weight,
and the precedent for creating it already exists. Mandated loan-level reporting
under the Home Mortgage Disclosure Act did not ask institutions to surrender
data they happened to hold; it required the record to be generated so that
oversight would have something to examine. A comparable and considerably
smaller requirement --- that a decision record bind to the operative model
version, configuration, and policy in force at decision time --- is the
smallest change that would convert several currently unreachable objects into
reachable ones. This paper does not propose that rule, and identifying the
smallest sufficient rule is not the same as establishing that it should be
adopted.

A third artifact would do more and cannot be produced at a desk: one worked
execution of the evidence request of Section~\ref{sec:request} at an actual
institution, recording which fields were answerable from records the bank held,
which required a compensating process, which were unanswerable because a vendor
controlled the evidence, and what the exercise cost in staff hours. That single
exhibit is what a regulator would need for a workability finding and what a
vendor would need for a demand signal.

\subsection{Theory of effect}

The paper assumes that clearer evidentiary distinctions can change what
architects preserve, what institutions claim, and what examiners ask. That
effect is plausible but unproven. Regulatory practice may instead stabilize
around easier documentary proxies; institutions may lack the authority or
resources to change vendor systems; or the vocabulary may be adopted without
the discipline it names. Coordination may produce another dashboard while
leaving underlying evidence inaccessible. Conversely, the frame may be useful even without
formal adoption if it helps one regulator identify an unsupported inference or
one architect preserve evidence that would otherwise disappear.

Adoption can fail by ceremonial compliance as well as rejection. Participants
might add ``affected party,'' ``binding,'' or ``indeterminacy'' fields to a
template while leaving the underlying artifact, access, and decision process
unchanged. Counts of completed fields would then measure use of the vocabulary,
not its effect. Evaluation should instead sample whether a claimed inference can
be traversed to its evidence, whether defeating observations are genuinely
available, whether declared scope reconciles to operational populations, and
whether identified gaps change a decision, disclosure, procurement request, or
policy analysis. Even those observations do not isolate causation without a
comparison, but they are closer to the proposed theory of effect than document
production alone.

The appropriate test is practical and revisable: do the diagnostics reveal
material inference gaps, do the proposed properties make those gaps more
inspectable, and do affected parties receive better representation as a result?
Evidence against any link should narrow or change the frame. This position paper
offers the questions and provisional capabilities; it does not claim the
outcomes of their adoption.

% Section 7: Conclusion

\section{Conclusion}
\label{sec:implications}

The FS~AI~RMF gives financial institutions a substantial control vocabulary
\parencite{treasuryFSAIRMFRelease2026}. It does not, by itself, tell a community
bank how those controls cohere or tell an examiner what a completed governance
artifact warrants. Empty-chair representation offers a proposed organizing
lens: identify the party absent from an AI-supported decision, the interest a
control is meant to represent, and the evidence by which that representation
can later be tested. The worked examples show how to apply the lens to selected
control objectives; framework-wide validation remains future work.

The paper's central distinction is between an artifact and the inference drawn
from it. Decision-time records, feature attributions, local approximations,
internal-state observations, generated rationales, human deliberative records,
and institutional approvals can all be useful. They expose different
evidentiary objects. An explanation reviewed by a human establishes that the
explanation was reviewed; stronger conclusions depend on what the method
targets, what the reviewer could test, and how the evidence is bound to the
claim. This distinction narrows neither AI nor explanation to a single approved
form. It requires institutions to state the proposition each artifact is being
used to support.

Silence-manufacture names a recurring inference gap: relevant evidence is
unavailable, a different artifact is offered, and the absence of contradictory
evidence supports a conclusion the artifact does not independently establish.
The pattern is a diagnostic, not a presumption of bad faith and not a claim that
every governance artifact fails. It directs attention to the relation among the
proposition, the artifact, the claimed binding, and the party who cannot test
that binding.

The three diagnostic techniques turn that distinction into a shared inquiry.
Empty-chair enumeration identifies whose interest depends on an
artifact. Pages missing from the ledger makes declared and ambient evidentiary
absences visible. The produced-silence question asks what access a design
removed and what is inferred from the resulting gap. Provisional evidentiary
capabilities then identify properties that may help participants answer those
questions. Vendors and architects are generally positioned to implement such
capabilities; banks can identify gaps, use compensating controls, and express
procurement needs; examiners can test claims within current authority; and
regulators can decide whether recurring gaps warrant coordinated expectations.
Their implementation, sufficiency, and tractability remain to be tested.

For a participant reading or contesting an AI governance record, the frame
reduces to five questions:

\begin{enumerate}
  \item What evidentiary object does this artifact expose?
  \item What conclusion is the institution asking the examiner to draw from it?
  \item What binds the artifact to that conclusion?
  \item Whose interest bears the cost if the binding is absent?
  \item What additional evidence would close or honestly disclose the gap?
\end{enumerate}

Those questions do not decide every governance dispute. They do make the point
of disagreement inspectable. A useful AI governance architecture should let an
absent party's interest appear not only in policy, but in evidence capable of
bearing the conclusion placed upon it.


\printbibliography

\end{document}
